% Source: sources/returns/S6_04_MAPS_MANIFOLDS/sections/manifolds.tex:684–713
\paragraph{Theorem D.1.11 (Known, with derivation: the diffeomorphism type of every cover).}
For each $d\geq1$ prime to six let $M_d$ denote the oriented
simply connected spin six-manifold with the following complete
Zhubr data: $A=\mathbb Z^{d-1}\oplus\mathbb Z/d$, $b_3=0$,
refined Pontryagin class zero, integral cubic zero, and finite
cyclic cubic as in D.1.9 in the basis
$x_0,\ldots,x_{d-2},u$. These data specify a unique oriented
diffeomorphism type. One has
\[
 \widetilde X_p\cong M_{|p|},\qquad M_1=S^6.
\]
For $d>1$ this is a classification by a complete published invariant;
no assertion that $M_d$ is a torsion-free connected sum is made.
\begin{proof}
We apply Zhubr's Theorem~1 \cite{manZhubr75}, with the definitions
on pp.~794--795. Its invariant quadruple consists of $H_2$, the
middle Betti number, the refined class $p$ and
$\mu=f_*[M]\in H_6(K(H_2,2);\mathbb Z)$, modulo isomorphism.
D.1.4 gives the first two data. D.1.8 gives $p_1=0$ and the refined
class satisfies $4p=\operatorname{PD}(p_1)$. Multiplication by four
is injective on $A$, so $p=0$. D.1.6 and D.1.9 give precisely all
the evaluations of $\mu$ enumerated in D.1.10. Thus the entire
quadruple is specified. Existence is supplied by the actual covers
already constructed, so its admissibility conditions require no
separate realization assumption. The injective classification in
Zhubr's theorem gives uniqueness and the displayed diffeomorphism.
For $d=1$ the same data are those of $S^6$. This also explains why
integral cup products alone, without D.1.9, would not have sufficed;
compare the torsion discussion in \cite{manReiser24} before Remark~4.7.
\end{proof}
